\NSPage{163}%
\begin{EESegment}{EE-TGT-P163-001}
Integration gives the correction's partial radial moment functions the same \NSi{NS-F-P163-001}{O_m(N^{-1})} bounds. Equation~\eqref{eq:C.16} therefore remains valid inside the correction interval, and
the admissible stress cone remains valid throughout that interval.
\end{EESegment}

\begin{EESegment}{EE-TGT-P163-002}
At the correction interval's right endpoint,
\NSi{NS-F-P163-002}{X_{\mathrm{rep}}}, the corrected profiles obey
\end{EESegment}
\NSFormula{NS-F-P163-003}{display}{C.17}
\begin{equation}
  \widetilde{\mathfrak m}(X_{\mathrm{rep}},\eta)
    =\mathfrak m(X_{\mathrm{rep}},\eta),
  \qquad
  (\widetilde E,\widetilde U)=(E,U)
  \quad\text{for }X\geq X_{\mathrm{rep}}.
  \tag{C.17}\label{eq:C.17}
\end{equation}
\begin{EESegment}{EE-TGT-P163-003}
From this endpoint onward, equality of the five radial moment functions continues by integration.
Lemma~\ref{lem:4.4} then gives equality of the complete pressure and the functions
\NSi{NS-F-P163-004}{Q_s,N_s} for every
\NSi{NS-F-P163-005}{X\geq X_{\mathrm{rep}}}, including equality of their parameter
derivatives. This keeps the terminal compensation already built into the input moment vector \NSi{NS-F-P163-006}{\mathfrak m}. As a result, the exterior moments
and the pressure normalized to vanish at radial infinity are unchanged.
The first collar is unchanged because both modifications are supported later, and the pressure is integrated from the unchanged axis datum. The last collar is
unchanged because the radial moments have been restored exactly. The two unused correction
patches still have their prescribed reference power law and
\NSi{NS-F-P163-007}{U=0}.
\end{EESegment}

\begin{EESegment}{EE-TGT-P163-004}
Finally, fix one such finite \NSi{NS-F-P163-008}{N}. Every mixed radial and
parameter derivative of the resulting profiles is finite, though its bound may depend on this \NSi{NS-F-P163-009}{N}. Fix this value before taking any limit in
the physical scale \NSi{NS-F-P163-010}{q}, and before choosing any later dyadic band or correction
stage.
\end{EESegment}
\end{proof}

\begin{EESegment}{EE-TGT-P163-005}
\subsection{The two edges and the completed profile.}\label{sec:C.3}

In this subsection, \NSi{NS-F-P163-011}{E,U,\Pi} mean the corrected
profiles from Proposition~\ref{prop:C.2}, and every related quantity is formed from those corrected profiles. Proposition~\ref{prop:C.2} has now supplied the interior cone inequalities and the exact radial moments.
It remains to prove Proposition~\ref{prop:C.3} by checking the remaining assertions at the annular
edges, where the stress \NSi{NS-F-P163-012}{\mathcal T_0} vanishes. Its direction
\NSi{NS-F-P163-013}{n} must extend smoothly, and the cone inequalities must keep a strict uniform margin. The derivatives of \NSi{NS-F-P163-014}{\mathcal T_0} must also obey the same
exponential weight \NSi{NS-F-P163-015}{\zeta} required in
Theorem~\ref{thm:4.6}. Both endpoint collars were left unchanged, so their original stress factorizations \eqref{eq:B.30} and
Equations~\eqref{eq:A.48} to~\eqref{eq:A.49} give these properties. In the weight below,
\NSi{NS-F-P163-016}{t_1>0} is the fixed activation width in
\NSi{NS-F-P163-017}{\log(X/X_a)} from Proposition~\ref{prop:4.10}.
\end{EESegment}

\begin{proposition}\label{prop:C.3}
\begin{EESegment}{EE-TGT-P163-006}
The profiles from Proposition~\ref{prop:C.2} have every property stated in
Theorem~\ref{thm:4.6}, with the weight
\end{EESegment}
\NSFormula{NS-F-P163-018}{display}{C.18}
\begin{equation}
\begin{gathered}
  y_a=\log(X/X_a),
  \qquad y_b=\log(X_b/X),
  \qquad \delta=\min\{1,y_a,y_b\},\\
  \zeta=\exp\!\left(-\frac{t_1^2}{y_a^2}-\frac{4}{y_b^2}\right)
  \qquad (X_a<X<X_b),
\end{gathered}
\tag{C.18}\label{eq:C.18}
\end{equation}
\begin{EESegment}{EE-TGT-P163-007}
and this weight is extended by zero when \NSi{NS-F-P163-019}{X\leq X_a} or
\NSi{NS-F-P163-020}{X\geq X_b}.
\end{EESegment}
\end{proposition}

\begin{proof}
\begin{EESegment}{EE-TGT-P163-008}
The stress is zero up to \NSi{NS-F-P163-021}{X_a} because of the analytic axis profile
and Proposition~\ref{prop:B.5}. It is zero beyond \NSi{NS-F-P163-022}{X_b} because of
Lemma~\ref{lem:A.8}. Proposition~\ref{prop:C.2} leaves both regions unchanged and makes
the admissible stress cone inequalities hold everywhere between them. If the
stress were zero at an interior point, the identity \eqref{eq:4.11} would force
\NSi{NS-F-P163-023}{p_s=(a,-b_s)}, and hence
\NSi{NS-F-P163-024}{P_{\mathrm c}=v_s}. That contradicts the strict quadratic cone test. This
proves the stress-support assertion in part~\textup{(ii)} of
Theorem~\ref{thm:4.6}. The unchanged axis profile also obeys \eqref{eq:4.13}
when \NSi{NS-F-P163-025}{X\leq X_a}.
\end{EESegment}

\begin{EESegment}{EE-TGT-P163-009}
For the cone assertions in part~\textup{(iii)}, set
\NSi{NS-F-P163-026}{n=\mathcal T_0/|\mathcal T_0|} on
\NSi{NS-F-P163-027}{X_a<X<X_b}. The interior bounds
\NSi{NS-F-P163-028}{F>0}, \NSi{NS-F-P163-029}{a>0}, and
\NSi{NS-F-P163-030}{v_s>2} come from the positive profile
construction and the admissible stress cone. At the inner endpoint,
\NSi{NS-F-P163-031}{a=p_{1,\mathrm r}>0}. Proposition~\ref{prop:B.5} gives
\NSi{NS-F-P163-032}{v_s>2+c} there, and the analytic axis profile has
\NSi{NS-F-P163-033}{F>0}. At the outer endpoint, the same bounds follow from
\eqref{eq:A.56}, together with \NSi{NS-F-P163-034}{b_s=0} and positivity of the heat
factor. Continuity on the compact closed annulus gives the required positive lower bounds throughout the annulus.
\end{EESegment}

\begin{EESegment}{EE-TGT-P163-010}
The inner edge keeps the factorization
\NSi{NS-F-P163-035}{\mathcal T_0=e_aB_0} from \eqref{eq:B.30}, where
\NSi{NS-F-P163-036}{e_a=e^{-t_1^2/y_a^2}g_a(y_a)},
\NSi{NS-F-P163-037}{g_a(0)>0}, and
\NSi{NS-F-P163-038}{B_0(0,\eta)=Fp_{s,\mathrm r}\neq0}. These facts make its direction
smooth.
\end{EESegment}

\NSPage{164}%
\begin{EESegment}{EE-TGT-P164-001}
At the edge, that direction is a
positive scalar multiple of the limiting shear
\NSi{NS-F-P164-001}{(a,-b_s)=a(1,t_s)}. In particular,
\NSi{NS-F-P164-002}{n_z-t_sn_\theta=0} and
\NSi{NS-F-P164-003}{n_\theta+t_sn_z>0} there. At the outer edge, the unchanged
factors from Proposition~\ref{prop:A.10} give
\end{EESegment}
\NSFormula{NS-F-P164-004}{display}{none}
\[
  n=\frac{(b_\theta,y_b^6b_z)}
          {\sqrt{b_\theta^2+y_b^{12}b_z^2}},
  \qquad b_\theta(0,\eta)>0.
\]
\begin{EESegment}{EE-TGT-P164-002}
The outer limiting direction is therefore \NSi{NS-F-P164-005}{(1,0)}, and the
limiting shear ratio is \NSi{NS-F-P164-006}{t_s=0}. Every statement here remains
smooth through \NSi{NS-F-P164-007}{\eta=\pm1}; the heat factor is used only where its argument is nonnegative.
\end{EESegment}

\begin{EESegment}{EE-TGT-P164-003}
These edge limits make the interior cone inequalities uniform on the closed annulus. In the interior, \eqref{eq:4.11} gives
\end{EESegment}
\NSFormula{NS-F-P164-008}{display}{none}
\[
  n_\theta+t_sn_z
    =\frac{P_{\mathrm c}-v_s}{|p_s-(a,-b_s)|}>0,
  \qquad
  n_z-t_sn_\theta
    =\frac{J_{\mathrm c}}{|p_s-(a,-b_s)|}.
\]
\begin{EESegment}{EE-TGT-P164-004}
The admissible stress cone makes
\NSi{NS-F-P164-009}{2-(v_s-2)(n_z-t_sn_\theta)^2/(n_\theta+t_sn_z)^2}
positive in the interior. This expression and its positive denominator both extend
continuously to the two edges. The factorizations just proved make the expression equal to two at either edge. The expression and the denominator each have a positive minimum on the compact annulus. Those two minima give one
\NSi{NS-F-P164-010}{0<\kappa<2} for which
\end{EESegment}
\NSFormula{NS-F-P164-011}{display}{none}
\[
  n_\theta+t_sn_z\geq\kappa,
  \qquad
  (v_s-2)(n_z-t_sn_\theta)^2
    \leq(2-\kappa)(n_\theta+t_sn_z)^2.
\]
\begin{EESegment}{EE-TGT-P164-005}
These bounds, together with the endpoint directions and the positive lower bounds above, prove part~\textup{(iii)} of Theorem~\ref{thm:4.6}.
\end{EESegment}

\begin{EESegment}{EE-TGT-P164-006}
For part~\textup{(iv)}, only the following weighted estimates remain to be proved.
\end{EESegment}
\NSFormula{NS-F-P164-012}{display}{C.19}
\begin{equation}
  |\mathcal T_0|\geq c\zeta,
  \qquad
  |\partial^I\mathcal T_0|\leq C_I\zeta\delta^{-m_I}
  \qquad (X_a<X<X_b,\ -1\leq\eta\leq1)
  \tag{C.19}\label{eq:C.19}
\end{equation}
\begin{EESegment}{EE-TGT-P164-007}
They must hold for every fixed mixed profile derivative
\NSi{NS-F-P164-013}{\partial^I}, with finite constants that do not depend on the physical
\NSi{NS-F-P164-014}{q}.
\end{EESegment}

\begin{EESegment}{EE-TGT-P164-008}
On an inner collar, \NSi{NS-F-P164-015}{\zeta} differs from
\NSi{NS-F-P164-016}{e^{-t_1^2/y_a^2}} only by a smooth positive factor bounded above
and below. The bounds in \eqref{eq:B.31} therefore give \eqref{eq:C.19} there. On
an outer collar, it differs from \NSi{NS-F-P164-017}{e^{-4/y_b^2}} in the same
way. The angular factor
\NSi{NS-F-P164-018}{e^{-4/y_b^2}y_b^{-3}b_\theta}, with
\NSi{NS-F-P164-019}{b_\theta} bounded below, gives the lower bound, while
\eqref{eq:A.51} gives the upper bound for every fixed derivative. On the compact interval left between the two collars, both \NSi{NS-F-P164-020}{\zeta} and
\NSi{NS-F-P164-021}{|\mathcal T_0|} have positive minima, and every
fixed derivative is bounded. Putting these three regions together proves
\eqref{eq:C.19}. The exponential factors also show that setting the stress equal to zero across both edges gives a smooth extension. The definition \eqref{eq:C.18} makes
\NSi{NS-F-P164-022}{\zeta} positive in the interior and flat at both edges. This completes part~\textup{(iv)} of Theorem~\ref{thm:4.6}.
\end{EESegment}

\begin{EESegment}{EE-TGT-P164-009}
For the regularity in part~\textup{(i)}, Proposition~\ref{prop:B.2} gives the
smooth axis profiles with \NSi{NS-F-P164-023}{F>0}, and
Corollary~\ref{cor:B.6} keeps the common analytic rectangle
\NSi{NS-F-P164-024}{[0,X_{\mathrm{an}}]\times[-1,1]}. The factors give Cartesian
smoothness across the axis through Equations~\eqref{eq:4.4} to~\eqref{eq:4.5}. The
later modifications in Proposition~\ref{prop:C.2} are smooth and preserve
positivity, with finite mixed derivative bounds on every fixed radial interval.
The inner directional margin is part~\textup{(iii)}, proved above. The heat
profile is smooth through \NSi{NS-F-P164-025}{\eta=\pm1} by
Lemma~\ref{lem:A.6}; its factor does not need to be analytic at zero. The pressure
normalization \eqref{eq:4.25} comes from Lemma~\ref{lem:A.5} and
Propositions~\ref{prop:A.7} and~\ref{prop:B.8}, and \eqref{eq:C.17} keeps it unchanged. This completes part~\textup{(i)}. Differentiating the pressure
formula gives \NSi{NS-F-P164-026}{\Pi_X=E^2/(2X)=F^2}, including its smooth extension
at \NSi{NS-F-P164-027}{X=0}. Proposition~\ref{prop:4.2} gives the exact tangential residual identities. Together with the support and axis checks above,
they complete part~\textup{(ii)}.
\end{EESegment}

\begin{EESegment}{EE-TGT-P164-010}
For part~\textup{(v)}, Propositions~\ref{prop:A.7} and~\ref{prop:B.8} give the exact exterior moments \eqref{eq:4.28}, and
\eqref{eq:C.17} keeps them unchanged. The same restoration keeps the heat exterior
\eqref{eq:4.29}; Lemma~\ref{lem:A.6} proves its exact heat evolution and smoothness. For the radius in this assertion, take
\NSi{NS-F-P164-028}{X_{\mathrm v}=X_{\mathrm{end}}\in(X_a,X_b)},
\end{EESegment}
